\chapter*{Conclusione}
\addcontentsline{toc}{chapter}{Conclusione}

Nel presente lavoro si è affrontato il problema del campionamento da distribuzioni di Boltzmann multimodali, al fine di stimare le quantità termodinamiche di un sistema a contatto termico con un serbatoio a temperatura ben definita. Si sono confrontate tre strategie concettualmente distinte: il campionamento markoviano locale (MCMC), i Normalizing Flows impiegati come generatori diretti con correzione tramite re-weighting, e lo schema ibrido Flow-MCMC. Il confronto è stato condotto su un sistema di test minimale ma non banale, un potenziale a doppio pozzo affiancato da direzioni armoniche.

I risultati dello studio al variare di dimensione, temperatura e numero di campioni mostrano un quadro coerente. MCMC locale, per costruzione vincolato a esplorare l'intorno immediato dello stato corrente, entra in crisi quando la barriera di potenziale diventa grande rispetto all'energia termica: il tasso di accettazione medio crolla al diminuire della temperatura e all'aumentare della dimensione del sistema, il tempo di autocorrelazione integrato $\tau^{int}$ della coordinata lenta cresce di ordini di grandezza e, nei casi più estremi, la catena rimane intrappolata in una sola delle due buche per l'intera simulazione.

Normalizing Flow e Flow-MCMC non soffrono di questa patologia. Flow-MCMC mantiene un tasso di accettazione medio prossimo a 1 su tutto l'intervallo di temperature esplorato; il tempo di autocorrelazione integrato $\tau^{int}$ e il suo analogo $\tau_{\text{eff}}$ restano vicini al valore ottimale, indipendentemente dalla temperatura. Coerentemente, entrambi i metodi convergono alle stime corrette con un numero di campioni sensibilmente inferiore rispetto a MCMC ed esplorano in modo bilanciato entrambi i pozzi a ogni temperatura considerata. In presenza di distribuzioni multimodali, si conclude dunque che le mosse globali apprese da un Normalizing Flow offrono un vantaggio pratico concreto rispetto al campionamento markoviano locale, sia come motore di campionamento diretto sia come proposta all'interno di uno schema Metropolis-Hastings.

I risultati ottenuti aprono naturalmente a diverse estensioni. Sarebbe auspicabile testare i metodi su potenziali che inducano distribuzioni di Boltzmann più complesse, così da verificare che la robustezza di NF e Flow-MCMC qui osservata si mantenga al crescere della complessità del panorama energetico. Altrettanto importante sarebbe studiare il costo computazionale dei tre metodi, finora confrontati solo in termini di efficienza statistica a parità di campioni generati. Infine, a mio giudizio, sarebbe di estremo interesse applicare direttamente lo schema Flow-MCMC a sistemi fisici realistici, caratterizzati da distribuzioni multimodali e da un numero elevato di gradi di libertà, come proteine, polimeri e materiali ferroelettrici e quantum-paraelettrici come gli ossidi perovskitici.